% =====================================================================
% Spec module — gemm_local_sram (spec baseline v1.0)
% Bố cục mục theo spec_1/modules/local_sram/latex/main.tex
% =====================================================================
\documentclass[10pt,a4paper]{article}
\newcommand{\DocHeader}{Spec module · gemm\_local\_sram}
\newcommand{\DocVersion}{baseline v1.0 · theo đặc tả hệ thống baseline v1.0}
\newcommand{\DocFooter}{Bộ tăng tốc GEMM INT8 trên Chipyard SoC — spec baseline}
\newcommand{\DocPdfTitle}{Spec module gemm_local_sram baseline v1.0}
% =====================================================================
% preamble.tex — gói, font, màu và định dạng chung (spec v1.0)
% Biên dịch bằng XeLaTeX. Bộ màu và cách đặt tiêu đề theo báo cáo tiến độ
% giai đoạn 1: đỏ = điều khiển/trạng thái, xanh = dữ liệu ma trận.
% =====================================================================
\usepackage[a4paper,top=2.2cm,bottom=2.2cm,left=2.1cm,right=2.1cm,headheight=15pt]{geometry}

% ---- Font: TeX Gyre (đủ dấu tiếng Việt, có sẵn trong TeX Live / Overleaf)
\usepackage{fontspec}
\setmainfont{texgyreheros}[
  Extension=.otf, UprightFont=*-regular, BoldFont=*-bold,
  ItalicFont=*-italic, BoldItalicFont=*-bolditalic]
\setsansfont{texgyreheros}[
  Extension=.otf, UprightFont=*-regular, BoldFont=*-bold,
  ItalicFont=*-italic, BoldItalicFont=*-bolditalic]
\newfontfamily\headfont{texgyreadventor}[
  Extension=.otf, UprightFont=*-regular, BoldFont=*-bold,
  ItalicFont=*-italic, BoldItalicFont=*-bolditalic]
\setmonofont{DejaVuSansMono}[
  Extension=.ttf, UprightFont=*, BoldFont=*-Bold,
  ItalicFont=*-Oblique, BoldItalicFont=*-BoldOblique, Scale=0.80]
% Khối code có chú thích tiếng Việt: DejaVu Sans Mono thiếu vài dấu, dùng TeX Gyre Cursor
\newfontfamily\vnmono{texgyrecursor}[
  NFSSFamily=vnmono, Extension=.otf, UprightFont=*-regular, BoldFont=*-bold,
  ItalicFont=*-italic, BoldItalicFont=*-bolditalic, Scale=0.9]

% ---- Ngôn ngữ
\usepackage{babel}
\babelprovide[import,main]{vietnamese}
\AtBeginDocument{%
  \renewcommand{\contentsname}{Mục lục}%
  \renewcommand{\listtablename}{Danh mục bảng}%
  \renewcommand{\listfigurename}{Danh mục hình}%
  \renewcommand{\tablename}{Bảng}%
  \renewcommand{\figurename}{Hình}}

% ---- Ký hiệu toán và Unicode
\usepackage{amsmath,amssymb}
\usepackage{newunicodechar}
\newunicodechar{→}{\ensuremath{\rightarrow}}
\newunicodechar{←}{\ensuremath{\leftarrow}}
\newunicodechar{↔}{\ensuremath{\leftrightarrow}}
\newunicodechar{⇒}{\ensuremath{\Rightarrow}}
\newunicodechar{≤}{\ensuremath{\le}}
\newunicodechar{≥}{\ensuremath{\ge}}
\newunicodechar{≠}{\ensuremath{\neq}}
\newunicodechar{≈}{\ensuremath{\approx}}
\newunicodechar{∈}{\ensuremath{\in}}
\newunicodechar{∉}{\ensuremath{\notin}}
\newunicodechar{∧}{\ensuremath{\wedge}}\newunicodechar{∨}{\ensuremath{\vee}}
\newunicodechar{⌈}{\ensuremath{\lceil}}
\newunicodechar{⌉}{\ensuremath{\rceil}}
\newunicodechar{Σ}{\ensuremath{\Sigma}}
\newunicodechar{−}{\ensuremath{-}}
\newunicodechar{×}{\ensuremath{\times}}
\newunicodechar{·}{\ensuremath{\cdot}}
\newunicodechar{²}{\ensuremath{^{2}}}
\newunicodechar{³}{\ensuremath{^{3}}}

% ---- Màu (khớp slides.html của báo cáo)
\usepackage{xcolor}
\definecolor{brand}{HTML}{BC1823}
\definecolor{brandsoft}{HTML}{F6DADC}
\definecolor{brandline}{HTML}{DF9DA2}
\definecolor{cream}{HTML}{FFF5EC}
\definecolor{card}{HTML}{FFFDF9}
\definecolor{ink}{HTML}{2B2B2B}
\definecolor{muted}{HTML}{6E5F5F}
\definecolor{data}{HTML}{1D5C7A}
\definecolor{datasoft}{HTML}{E2EDF2}
\definecolor{cpu}{HTML}{4A4040}
\definecolor{okgreen}{HTML}{2E7D4F}
\color{ink}

% ---- Tiêu đề mục: chữ đỏ, vạch đỏ bên trái như tiêu đề slide
\usepackage{titlesec}
\titleformat{\section}[block]
  {\headfont\bfseries\LARGE\color{brand}}
  {\thesection}{0.6em}{}
  [{\color{brandline}\titlerule[0.8pt]}]
\titlespacing*{\section}{0pt}{14pt}{8pt}
\titleformat{\subsection}[block]
  {\headfont\bfseries\large\color{brand}}
  {\thesubsection}{0.6em}{}
\titlespacing*{\subsection}{0pt}{12pt}{4pt}
\titleformat{\subsubsection}[block]
  {\headfont\bfseries\normalsize\color{ink}}
  {\thesubsubsection}{0.6em}{}
\titlespacing*{\subsubsection}{0pt}{8pt}{2pt}
\setcounter{secnumdepth}{3}
\setcounter{tocdepth}{2}

% ---- Bảng, danh sách, code
\usepackage{array,longtable,booktabs,multirow}
\usepackage{ragged2e}
\usepackage{enumitem}
\setlist{itemsep=1pt,topsep=3pt,leftmargin=1.4em}
\usepackage{fancyvrb}
\fvset{fontsize=\small,fontfamily=vnmono}
\setlength{\LTpre}{4pt}
\setlength{\LTpost}{6pt}
\setlength{\LTcapwidth}{\linewidth}
\renewcommand{\arraystretch}{1.18}
\newcolumntype{L}[1]{>{\RaggedRight\arraybackslash}p{\dimexpr#1\linewidth-2\tabcolsep-0.6pt\relax}}
\newcolumntype{C}[1]{>{\Centering\arraybackslash}p{\dimexpr#1\linewidth-2\tabcolsep-0.6pt\relax}}

% Bảng chuẩn: \begin{spectable}{cột}{chú thích}{nhãn}{hàng tiêu đề} ... \end{spectable}
% Hàng tiêu đề lặp lại khi bảng sang trang.
\newenvironment{spectable}[4]{%
  \begingroup\small
  \begin{longtable}{@{}#1@{}}
  \caption{#2}\label{#3}\\
  \toprule #4 \\ \midrule
  \endfirsthead
  \caption[]{#2 (tiếp theo)}\\
  \toprule #4 \\ \midrule
  \endhead
  \bottomrule
  \endlastfoot}
  {\end{longtable}\endgroup}
% Bảng yêu cầu ID | Yêu cầu
\newenvironment{reqtable}[2]{%
  \begin{spectable}{L{0.12}L{0.88}}{#1}{#2}%
  {\textbf{ID} & \textbf{Yêu cầu}}}
  {\end{spectable}}

% ---- Chú thích, đánh số theo mục
\usepackage{caption}
\captionsetup{font=small,labelfont={bf,color=brand},skip=4pt}
\counterwithin{table}{section}
\counterwithin{figure}{section}

% ---- Hình (TikZ)
\usepackage{graphicx}
\usepackage{tikz}
\usetikzlibrary{arrows.meta,positioning,fit,calc,backgrounds,shapes.geometric,decorations.pathreplacing}
\tikzset{
  >={Stealth[length=2.2mm]},
  blk/.style={draw=brand,line width=0.8pt,fill=card,rounded corners=2pt,align=center,font=\small,inner sep=4pt},
  blkd/.style={draw=data,line width=0.8pt,fill=datasoft,rounded corners=2pt,align=center,font=\small,inner sep=4pt},
  blkn/.style={draw=brand,line width=1.1pt,fill=brandsoft,rounded corners=2pt,align=center,font=\small,inner sep=4pt},
  blkc/.style={draw=cpu,line width=0.8pt,fill=white,rounded corners=2pt,align=center,font=\small,inner sep=4pt},
  ctl/.style={->,draw=brand,line width=0.8pt},
  dat/.style={->,draw=data,line width=1.3pt},
  lbl/.style={font=\scriptsize,align=center,inner sep=1pt},
}
\usepackage{float}
\usepackage{placeins}
\usepackage{needspace}
\AddToHook{cmd/section/before}{\FloatBarrier\Needspace{12\baselineskip}}
\AddToHook{cmd/subsection/before}{\Needspace{10\baselineskip}}
\AddToHook{cmd/subsection/before}{\FloatBarrier}

% ---- Hộp nội dung kiểu báo cáo
\usepackage[most]{tcolorbox}
% Hộp kết luận đầu mục (một câu trả lời trước, chi tiết sau)
\newtcolorbox{keypoint}{enhanced,breakable,colback=brandsoft!55,colframe=brand,
  boxrule=0pt,leftrule=3pt,arc=0pt,left=8pt,right=6pt,top=4pt,bottom=4pt,
  fontupper=\small}
% Hộp ghi chú
\newtcolorbox{notebox}[1][Ghi chú]{enhanced,breakable,colback=datasoft!60,colframe=data,
  boxrule=0pt,leftrule=3pt,arc=0pt,left=8pt,right=6pt,top=4pt,bottom=4pt,
  fontupper=\small,title={\headfont\bfseries\footnotesize\MakeUppercase{#1}},
  coltitle=data,colbacktitle=datasoft!60,attach title to upper={\par}}
% Bốn câu hỏi của mỗi giao tiếp (như slide "Giao tiếp 1/4")
\newcommand{\qitem}[2]{{\headfont\bfseries\footnotesize\color{brand}\MakeUppercase{#1}}\par\vspace{1pt}#2\par\vspace{5pt}}
\newtcolorbox{qcard}[1]{enhanced,breakable,colback=card,colframe=brandline,boxrule=0.6pt,
  arc=2pt,left=7pt,right=7pt,top=5pt,bottom=3pt,fontupper=\small,
  title={#1},fonttitle=\headfont\bfseries\small,coltitle=brand,colbacktitle=card,
  titlerule=0.4pt,titlerule style={brandline}}
\newcommand{\interfacecard}[5]{%
  \begin{qcard}{#1}
  \qitem{Trao đổi gì?}{#2}
  \qitem{Ai khởi tạo?}{#3}
  \qitem{Khi nào được chấp nhận?}{#4}
  \qitem{Khi nào sang bước tiếp?}{#5}
  \end{qcard}}

% Nhãn trạng thái giai đoạn
\newcommand{\stagetag}[1]{\tikz[baseline=(t.base)]\node[t,fill=brand,text=white,rounded corners=1.5pt,inner sep=1.6pt,font=\headfont\bfseries\scriptsize](t){#1};}
\tikzset{t/.style={}}

% ---- Đầu trang, chân trang
\usepackage{fancyhdr}
\pagestyle{fancy}
\fancyhf{}
\fancyhead[L]{\headfont\footnotesize\color{brand}\MakeUppercase{Spec module · \ModHeader}}
\fancyhead[R]{\headfont\footnotesize\color{muted}v1.0 · spec hệ thống v1.0, gồm mục 7}
\fancyfoot[L]{\footnotesize\color{muted}Bộ tăng tốc GEMM INT8 trên Chipyard SoC}
\fancyfoot[R]{\headfont\bfseries\footnotesize\color{brand}\thepage}
\renewcommand{\headrule}{{\color{brandline}\hrule height 0.6pt}}
\renewcommand{\footrule}{{\color{brandline}\hrule height 0.6pt}}
\renewcommand{\footrulewidth}{0.6pt}
\fancypagestyle{plain}{\fancyhf{}\renewcommand{\headrulewidth}{0pt}%
  \fancyfoot[R]{\headfont\bfseries\footnotesize\color{brand}\thepage}}

\setlength{\parindent}{0pt}
\setlength{\parskip}{4pt}
\linespread{1.08}

% ---- Lệnh riêng
\newcommand{\XN}{XN}
\newcommand{\appsection}[1]{%
  \refstepcounter{section}%
  \section*{Phụ lục \thesection: #1}%
  \addcontentsline{toc}{section}{Phụ lục \thesection: #1}}

\usepackage[hidelinks]{hyperref}
\usepackage{url}
\DeclareUrlCommand\sig{\urlstyle{tt}}
\makeatletter\g@addto@macro\UrlBreaks{\do\_\do\[}\makeatother
% Bản dùng được trong chú thích và tiêu đề (đối số di chuyển)
\DeclareRobustCommand\sigc[1]{\texttt{\detokenize{#1}}}
\hypersetup{pdftitle={Spec module gemm_systolic_array v1.0},pdfauthor={Nhóm đồ án GEMM}}

\newcommand{\eqbox}[1]{\begin{tcolorbox}[enhanced,colback=card,colframe=brandline,boxrule=0.6pt,arc=2pt,left=8pt,right=8pt,top=2pt,bottom=2pt]#1\end{tcolorbox}}

% ---- Bìa dùng chung: \makecover{tên module}{tiêu đề phụ}{dòng mô tả}{các dòng bảng}
\newcommand{\makecover}[4]{%
\thispagestyle{empty}
\begin{tikzpicture}[remember picture,overlay]
  \fill[cream] (current page.north west) rectangle (current page.south east);
\end{tikzpicture}
\vspace*{3cm}
{\headfont\bfseries\small\color{brand}\MakeUppercase{Spec module · phiên bản 1.0 · full option}\par}
\vspace{0.6cm}
\begin{tikzpicture}[baseline=(b.north)]\node[fill=brand,minimum width=0.25cm,minimum height=3.6cm,inner sep=0pt](b){};\end{tikzpicture}\hspace{0.45cm}\begin{minipage}[t]{0.9\linewidth}
{\headfont\bfseries\fontsize{28}{34}\selectfont\color{brand}#1\par}
\vspace{6pt}
{\headfont\bfseries\fontsize{18}{24}\selectfont\color{brand}#2\par}
\vspace{10pt}
{\large\color{ink}#3\par}
\vspace{4pt}
{\normalsize\color{muted}Bộ tăng tốc GEMM INT8 trên Chipyard SoC\par}
\end{minipage}
\vspace{2cm}

\begingroup\small
\begin{tabular}{@{}L{0.2}L{0.78}@{}}
\toprule
#4
\bottomrule
\end{tabular}
\endgroup
\clearpage
\tableofcontents
\clearpage}


\definecolor{brand}{HTML}{000000}
\definecolor{brandline}{HTML}{888888}
\definecolor{ink}{HTML}{000000}
\color{black}
\titleformat{\section}{\bfseries\Large}{\thesection.}{0.6em}{}
\titleformat{\subsection}{\bfseries\large}{\thesubsection.}{0.6em}{}
\titleformat{\subsubsection}{\bfseries}{\thesubsubsection.}{0.6em}{}
\counterwithout{table}{section}
\counterwithout{figure}{section}
\captionsetup{labelfont=bf,justification=raggedright,singlelinecheck=false}
\fancyhead[L]{\small Đặc tả thiết kế GEMM INT8}
\fancyhead[R]{\small\ModHeader}
\fancyfoot[L]{\small Bộ tăng tốc GEMM INT8}
\fancyfoot[R]{\small\thepage}
\hypersetup{pdftitle={Đặc tả thiết kế GEMM INT8},pdfauthor={Nhóm đồ án GEMM}}

\titleformat{\section}{\color{data}\bfseries\Large}{\thesection.}{0.6em}{}
\titleformat{\subsection}{\color{data}\large}{\thesubsection.}{0.6em}{}
\renewenvironment{spectable}[4]{%
  \begingroup\small
  \begin{longtable}{#1}
  \caption{#2}\label{#3}\\
  \hline #4 \\ \hline
  \endfirsthead
  \caption[]{#2 (tiếp theo)}\\
  \hline #4 \\ \hline
  \endhead
  \endlastfoot}
  {\end{longtable}\endgroup}
\setlength{\emergencystretch}{2em}

\begin{document}
{\Large\bfseries ĐẶC TẢ THIẾT KẾ gemm\_local\_sram}\par
\vspace{2pt}
{\small\color{muted}Module của \sig{gemm_accel_top} · instance \sig{u_sram}}\par
\tableofcontents
\clearpage

% ===================== 1 =====================
\section{Tổng quan}
\subsection{Tính năng chính}\label{sec:1}

Scratchpad giữ bốn buffer logic A, B, C, BIAS, mỗi buffer có hai bank (ping-pong, NB = 2). Phía tính (array, VPU) cần đọc hoặc ghi 8 phần tử trong một chu kỳ, nên mỗi bank được chia thành nhiều \textbf{sub-bank} một cổng sao cho mỗi truy cập song song chạm các sub-bank khác nhau. Phía DMA đọc hoặc ghi một word mỗi chu kỳ. Hai phía dùng hai bank khác nhau, nên không bao giờ tranh nhau một sub-bank.

\paragraph{Chức năng}
\begin{itemize}
\item Mỗi buffer có hai bank. Mỗi bank: A 256 word, B 256 word, C 1024 word, BIAS 32 word, tức 1568 word mỗi bank; tổng 3136 word = 12,25 KiB
\item Mọi cổng đọc có độ trễ cố định đúng 1 chu kỳ; mọi lệnh ghi có hiệu lực ở cạnh clock cuối chu kỳ có lệnh; không có backpressure (\sig{o_dw_ready} = \sig{o_dr_ready} = 1 hằng số)
\item Mỗi sub-bank là SRAM một cổng đồng bộ. Bank b của buffer X chỉ bị truy cập bởi một phía trong một chu kỳ; hai bank khác nhau được truy cập độc lập cùng chu kỳ
\item Mỗi chu kỳ phía tính truy cập đồng thời A, B, BIAS, C; phía DMA truy cập đồng thời một cổng ghi (A, B hoặc BIAS) và một cổng đọc (C)
\item Nội dung RAM không được khởi tạo bởi \sig{rst_n} hay \sig{soft_clr}; nội dung sau reset không xác định
\end{itemize}

% ===================== 2 =====================
\subsection{Tham số}\label{sec:2}
\begin{spectable}{|L{0.2}|L{0.14}|L{0.66}|}{Tham số}{tab:2-1}{\textbf{Tham số} & \textbf{Giá trị} & \textbf{Ý nghĩa}}
\sig{P} & 8 & localparam cố định, bằng \sig{PE_ARRAY_SIZE} của top (chỉ nhận 8): số lane của vector phía tính \\\hline

\sig{NB}, \sig{Q}, \sig{LP} & 2, 2, 3 & localparam. NB = 2 luôn luôn (hai bank mỗi buffer). Q = P/4 = 2 (số sub-bank của B); LP = $\log_2 P$ = 3 \\\hline

\end{spectable}

\begin{notebox}[Quyết định baseline]
\textbf{D-04 (P = 8 cố định).} Q = P/4 = 2, LP = 3; bus vector A, B rộng 8P = 64 bit, bus vector BIAS và C rộng 32P = 256 bit. Mọi công thức trong tài liệu này được viết cho P = 8 và không có biến thể P = 4.
\end{notebox}

% ===================== 3 =====================
\subsection{Tín hiệu giao tiếp}\label{sec:3}

\begin{spectable}{|L{0.26}|L{0.15}|L{0.14}|L{0.45}|}{Tín hiệu giao tiếp của \sigc{gemm_local_sram}}{tab:3-1}{\textbf{Tên tín hiệu} & \textbf{Độ rộng} & \textbf{Hướng} & \textbf{Mô tả}}
\sig{clk} & 1 & Input & Clock của khối.\newline Các thanh ghi cập nhật tại cạnh lên. \\\hline
\sig{rst_n} & 1 & Input & Reset tích cực mức thấp; được kích hoạt bất đồng bộ.\newline Cạnh nhả reset được đồng bộ với clock tại top trước khi đưa tới các module. \\\hline
\sig{i_soft_clr} & 1 & Input & Xung xóa đồng bộ dài một chu kỳ do Main FSM phát.\newline Xóa trạng thái điều khiển tại cạnh lên; phạm vi xóa dữ liệu được quy định trong mục Reset. \\\hline


\multicolumn{4}{@{}l}{\textit{Cổng ghi phía DMA (LOAD\_A, LOAD\_B, LOAD\_BIAS)}} \\\hline

\sig{i_dw_valid} & 1 & Input & Yêu cầu ghi một word 32 bit vào scratchpad.\newline Ghi được nhận khi valid và ready cùng bằng 1 tại cạnh lên. \\\hline

\sig{i_dw_buf} & 2 & Input & Chọn buffer đích cho word ghi.\newline Giá trị 0: A.\newline Giá trị 1: B.\newline Giá trị 2: BIAS; giá trị 3 không hợp lệ. \\\hline

\sig{i_dw_bank} & 1 & Input & Chọn bank đích cho word ghi.\newline Hợp lệ cùng yêu cầu ghi; nhận giá trị 0 hoặc 1. \\\hline

\sig{i_dw_addr} & 10 & Input & Địa chỉ word logic trong buffer đích.\newline A/B dùng miền 0 đến 255; BIAS dùng miền 0 đến 31.\newline Hợp lệ cùng yêu cầu ghi. \\\hline

\sig{i_dw_data} & 32 & Input & Word dữ liệu 32 bit cần ghi vào scratchpad.\newline Hợp lệ cùng yêu cầu ghi; DMA giữ nguyên thứ tự byte của dữ liệu DRAM. \\\hline

\sig{o_dw_ready} & 1 & Output & Báo scratchpad nhận word ghi.\newline Luôn bằng 1 (không có backpressure); giao dịch vẫn được xác định bởi valid và ready. \\\hline


\multicolumn{4}{@{}l}{\textit{Cổng đọc phía DMA (STORE\_C)}} \\\hline

\sig{i_dr_valid} & 1 & Input & Yêu cầu đọc một word từ buffer C.\newline Yêu cầu được nhận tại cạnh lên khi valid và ready cùng bằng 1. \\\hline
\sig{i_dr_bank} & 1 & Input & Chọn bank nguồn của buffer C.\newline Hợp lệ cùng yêu cầu đọc; nhận giá trị 0 hoặc 1. \\\hline
\sig{i_dr_addr} & 10 & Input & Địa chỉ word logic cần đọc trong buffer C.\newline Hợp lệ cùng yêu cầu đọc; miền địa chỉ phụ thuộc bố cục INT32 hoặc INT8 của C. \\\hline

\sig{o_dr_ready} & 1 & Output & Báo scratchpad nhận yêu cầu đọc C.\newline Luôn bằng 1 (không có backpressure). \\\hline

\sig{o_dr_rvalid} & 1 & Output & Báo word đọc trả về hợp lệ.\newline Xuất hiện đúng một chu kỳ sau yêu cầu được nhận; mỗi yêu cầu có đúng một phản hồi, theo thứ tự. \\\hline
\sig{o_dr_rdata} & 32 & Output & Word dữ liệu 32 bit đọc từ buffer C.\newline Bên nhận lấy dữ liệu khi read rvalid = 1; không có ready cho kênh trả dữ liệu. \\\hline


\multicolumn{4}{@{}l}{\textit{Bank phía tính (từ FSM)}} \\\hline

\sig{i_comp_bank_a} & 1 & Input & Chọn bank của buffer A dành cho phía tính.\newline Giữ nguyên trong suốt COMPUTE. \\\hline
\sig{i_comp_bank_b} & 1 & Input & Chọn bank của buffer B dành cho phía tính.\newline Giữ nguyên trong suốt COMPUTE. \\\hline
\sig{i_comp_bank_bias} & 1 & Input & Chọn bank của buffer BIAS dành cho phía tính.\newline Giữ nguyên trong suốt COMPUTE. \\\hline
\sig{i_comp_bank_c} & 1 & Input & Chọn bank của buffer C dành cho phía tính.\newline Giữ nguyên trong suốt COMPUTE. \\\hline

\sig{i_comp_active} & 1 & Input & Báo COMPUTE đang hoạt động.\newline SRAM dùng tín hiệu này để kiểm tra quyền truy cập bằng assertion; tín hiệu không thay thế các enable đọc/ghi. \\\hline


\multicolumn{4}{@{}l}{\textit{Cổng đọc A, B (từ array)}} \\\hline

\sig{i_a_rd_en} & 1 & Input & Cho phép đọc một vector A.\newline Địa chỉ hợp lệ trong chu kỳ enable = 1; dữ liệu trả về đúng một chu kỳ sau. \\\hline
\sig{i_a_rd_row} & 5 & Input & Hàng đầu i0 của vector A.\newline Là bội của 8; hợp lệ cùng read enable. \\\hline
\sig{i_a_rd_k} & 5 & Input & Chỉ số k của cột A cần đọc.\newline Hợp lệ cùng read enable; phải nhỏ hơn K. \\\hline

\sig{o_a_rd_data} & 64 & Output & Vector 8 byte INT8 của A.\newline Lane r tại bit [8r +: 8] chứa A[i0+r][k].\newline Hợp lệ một chu kỳ sau yêu cầu đọc; lane ngoài miền tác vụ chứa giá trị bất kỳ. \\\hline

\sig{i_b_rd_en} & 1 & Input & Cho phép đọc một vector B.\newline Địa chỉ hợp lệ trong chu kỳ enable = 1; dữ liệu trả về đúng một chu kỳ sau. \\\hline
\sig{i_b_rd_k} & 5 & Input & Chỉ số k của hàng B cần đọc.\newline Hợp lệ cùng read enable; phải nhỏ hơn K. \\\hline
\sig{i_b_rd_col} & 5 & Input & Cột đầu j0 của vector B.\newline Là bội của 8; hợp lệ cùng read enable. \\\hline

\sig{o_b_rd_data} & 64 & Output & Vector 8 byte INT8 của B.\newline Lane c tại bit [8c +: 8] chứa B[k][j0+c].\newline Hợp lệ một chu kỳ sau yêu cầu đọc; lane ngoài miền tác vụ chứa giá trị bất kỳ. \\\hline


\multicolumn{4}{@{}l}{\textit{Cổng BIAS và C (từ VPU)}} \\\hline

\sig{i_bias_rd_en} & 1 & Input & Yêu cầu đọc 8 phần tử bias.\newline Bằng 1 khi đoạn hàng hợp lệ và hậu xử lý được bật; dữ liệu về sau một chu kỳ. \\\hline
\sig{i_bias_rd_col} & 5 & Input & Cột đầu j0 của vector bias.\newline Là bội của 8; hợp lệ cùng bias read enable. \\\hline

\sig{o_bias_rd_data} & 256 & Output & Vector 8 giá trị bias INT32.\newline Lane c tại bit [32c +: 32] chứa bias[j0+c].\newline Hợp lệ đúng một chu kỳ sau yêu cầu; lane bị mask chứa giá trị bất kỳ. \\\hline

\sig{i_c_wr_en} & 1 & Input & Cho phép ghi một đoạn hàng kết quả vào buffer C tại cạnh lên.\newline Các trường fmt, row, col, data và mask phải hợp lệ trong chu kỳ này. \\\hline
\sig{i_c_wr_fmt} & 1 & Input & Chọn định dạng phần tử ghi C.\newline Giá trị 0: INT32, mỗi phần tử là một word.\newline Giá trị 1: INT8, mỗi phần tử là một byte.\newline Hợp lệ cùng write enable. \\\hline
\sig{i_c_wr_row} & 5 & Input & Chỉ số hàng i của buffer C cần ghi.\newline Hợp lệ cùng write enable. \\\hline
\sig{i_c_wr_col} & 5 & Input & Cột đầu j0 của đoạn hàng cần ghi.\newline Là bội của 8; hợp lệ cùng write enable. \\\hline

\sig{i_c_wr_data} & 256 & Input & Dữ liệu ghi đoạn hàng C.\newline Với fmt = 0, phần tử c nằm tại bit [32c +: 32].\newline Với fmt = 1, phần tử c nằm tại bit [8c +: 8] và phần cao còn lại bằng 0.\newline Hợp lệ cùng write enable. \\\hline
\sig{i_c_wr_mask} & 8 & Input & Cho phép ghi theo phần tử.\newline Bit c = 1 ghi phần tử c theo định dạng fmt; bit c = 0 giữ nguyên phần tử đó.\newline Hợp lệ cùng write enable. \\\hline

\end{spectable}

\paragraph{Giao tiếp}
\begin{itemize}
\item \sig{o_dw_ready} và \sig{o_dr_ready} luôn bằng 1. Lệnh chỉ được nhận khi valid và ready cùng bằng 1. Mỗi chu kỳ nhận một word ghi và một yêu cầu đọc; mỗi yêu cầu đọc có đúng một phản hồi theo thứ tự
\item \sig{o_a_rd_data}, \sig{o_b_rd_data}, \sig{o_bias_rd_data} hợp lệ ở đúng chu kỳ sau chu kỳ có \sig{*_rd_en}; ở chu kỳ khác giá trị tùy ý. \sig{o_dr_rvalid} là thanh ghi
\item Lane ứng với hàng hoặc cột ngoài vùng hợp lệ của tác vụ mang nội dung tùy ý
\end{itemize}

\subsection{Điều kiện phục vụ đồng thời}
Trong một chu kỳ SRAM phải phục vụ đồng thời một lần đọc A (8 byte từ 8 hàng), một lần đọc B (8 byte liền nhau), một lần đọc BIAS (8 word) và một lần ghi C (8 phần tử có mask). Hai phía tính và DMA chạy đồng thời trên các bank khác nhau.

\begin{spectable}{|L{0.2}|L{0.2}|L{0.6}|}{Truy cập song song và sub-bank bị chạm (P = 8)}{tab:3-2}{\textbf{Truy cập} & \textbf{Sub-bank chạm} & \textbf{Vì sao không trùng}}
Đọc A (\sig{i_a_rd_*}) & 8 của A[b] & $i_0$ là bội của 8, nên 8 hàng $i_0 + r$ có $i \bmod 8$ khác nhau \\\hline
Đọc B (\sig{i_b_rd_*}) & 2 của B[b] & $j_0$ là bội của 8, nên hai word $jw = j_0/4,\ j_0/4+1$ có $jw \bmod 2$ khác nhau \\\hline
Đọc BIAS (\sig{i_bias_rd_*}) & 8 của BIAS[b] & 8 cột $j_0 + c$ có $j \bmod 8$ khác nhau \\\hline
Ghi C INT32 (\sig{i_c_wr_*}, fmt = 0) & 8 của C[b] & 8 cột $j_0 + c$ có $j \bmod 8$ khác nhau \\\hline
Ghi C INT8 (fmt = 1) & 2 của C[b] & hai word liền nhau $jw$ có $jw \bmod 8$ khác nhau \\\hline
DMA ghi A, B, BIAS & 1 & một word mỗi chu kỳ \\\hline
DMA đọc C & 1 & một word mỗi chu kỳ \\\hline

\end{spectable}

\begin{itemize}
\item Phía tính chạm bank \sig{i_comp_bank_*} của cả bốn buffer; phía DMA chạm buffer và bank do \sig{i_dw_*}, \sig{i_dr_*} chọn. Các buffer khác nhau là các RAM khác nhau nên không tranh nhau; trong cùng một buffer, Main FSM cấp hai bank khác nhau cho hai phía (scoreboard) khi phía tính thực sự dùng buffer đó, được kiểm bằng AS-02. Riêng BIAS: job có \sig{POST_EN} = 0 vẫn mang một giá trị \sig{i_comp_bank_bias} nhưng không đọc BIAS, nên DMA được phép ghi BIAS vào bank đó
\item Đọc A/B/BIAS trả dữ liệu ở chu kỳ kế tiếp; bên nhận chỉ dùng các lane hợp lệ. Ghi C với fmt = 0 lấy phần tử c từ \sigc{data[32c +: 32]} và \sig{mask[c]} cho phép ghi cả word. Với fmt = 1, lấy byte c từ \sigc{data[8c +: 8]}, \sig{mask[c]} cho phép ghi byte đó; phần cao \sig{data[255:64]} bằng 0. Hàng/cột đầu phải hợp lệ và cột đầu là bội của 8. Mask chặn các phần tử ngoài miền tác vụ
\item Mỗi yêu cầu đọc DMA được nhận phải tạo đúng một \sig{o_dr_rvalid} ở chu kỳ tiếp theo, theo đúng thứ tự; không có ready cho dữ liệu trả về. Nội dung RAM và các lane không hợp lệ không cần khởi tạo, và được phép chứa X trong mô phỏng
\end{itemize}

% ===================== 4 =====================
\section{Đặc tả thanh ghi}
Module không có thanh ghi ánh xạ MMIO. Toàn bộ thanh ghi MMIO nằm trong \sig{gemm_register_file}, xem đặc tả hệ thống baseline. Cấu hình được nhận qua các tín hiệu đầu vào (bank phía tính từ Main FSM, lệnh từ DMA, array và VPU); trạng thái nội bộ, điều kiện chốt và reset được mô tả bên dưới.

\section{Mô tả chức năng}
\subsection{Bố cục logic}\label{sec:4}

\begin{spectable}{|L{0.12}|L{0.25}|L{0.63}|}{Địa chỉ word logic trong mỗi bank}{tab:4-1}{\textbf{Buffer} & \textbf{Địa chỉ a} & \textbf{Nội dung word}}
A & \{i[4:0], kw[2:0]\}, kw = k » 2 & Byte e = $A[i][4kw + e]$; lane = k[1:0] \\\hline

B & \{k[4:0], jw[2:0]\}, jw = j » 2 & Byte e = $B[k][4jw + e]$ \\\hline

BIAS & \{j[4:0]\} & $bias[j]$ \\\hline

C, INT32 & \{i[4:0], j[4:0]\} & $C[i][j]$ (word 0 … 1023) \\\hline

C, INT8 & \{i[4:0], jw[2:0]\} & Byte e = $Y[i][4jw + e]$, lane = j[1:0] (dùng word 0 … 255) \\\hline

\end{spectable}

\begin{itemize}
\item Mỗi phần tử A, B là INT8 có dấu; bốn phần tử liên tiếp theo K (A) hoặc theo N (B) nằm trong một word 32 bit, phần tử có chỉ số thấp ở byte thấp (little-endian, trùng thứ tự byte của DRAM)
\item Cùng một buffer C chứa INT32 (khi \sig{POST_EN} = 0) hoặc INT8 (khi \sig{POST_EN} = 1); định dạng do \sig{i_c_wr_fmt} chọn ở mỗi lệnh ghi và do \sig{i_wr_fmt} của DMA chọn khi đọc. SRAM không lưu định dạng
\item Quy tắc đóng gói INT8 phía VPU: phần tử c nằm ở \sig{i_c_wr_data[8c+7:8c]}, \sig{i_c_wr_data[255:64]} = 0; \sig{i_c_wr_mask[c]} là byte enable của đúng byte đó
\end{itemize}

% ===================== 5 =====================
\subsection{Tổ chức vật lý và ánh xạ}\label{sec:5}

\begin{spectable}{|L{0.1}|L{0.12}|L{0.2}|L{0.2}|L{0.38}|}{Sub-bank của mỗi bank (P = 8, Q = 2, LP = 3)}{tab:5-1}{\textbf{Buffer} & \textbf{Số sub-bank} & \textbf{Sâu × rộng} & \textbf{Sub-bank của a} & \textbf{Địa chỉ cục bộ}}
A & 8 & 32 × 32 bit & i mod 8 = i[2:0] & \{i[4:3], kw\} (5 bit) \\\hline

B & Q = 2 & 128 × 32 bit & jw mod 2 = jw[0] & \{k, jw » 1\} (7 bit) \\\hline

BIAS & 8 & 4 × 32 bit, register file & j mod 8 = j[2:0] & j » 3 = j[4:3] (2 bit) \\\hline

C & 8 & 128 × 32 bit, ghi theo byte & a mod 8 = a[2:0] & a » 3 = a[9:3] (7 bit) \\\hline

\end{spectable}

\begin{spectable}{|L{0.22}|L{0.2}|L{0.14}|L{0.14}|L{0.3}|}{Kiểm kê bộ nhớ vật lý (NB = 2 bank)}{tab:5-2}{\textbf{Bộ nhớ} & \textbf{Kích thước một RAM} & \textbf{Số word} & \textbf{Số RAM} & \textbf{Ghi chú}}
\sig{spm_a_subbank} & 1024 bit & 32 & 16 & 8 sub-bank × 2 bank; một cổng đọc/ghi \\\hline
\sig{spm_b_subbank} & 4096 bit & 128 & 4 & 2 sub-bank × 2 bank \\\hline
\sig{spm_bias_subbank_reg} & 128 bit & 4 & 16 & 8 sub-bank × 2 bank; register file \\\hline
\sig{spm_c_subbank_be} & 4096 bit & 128 & 16 & 8 sub-bank × 2 bank; byte enable \\\hline

\end{spectable}

Tổng dung lượng: 16 × 1024 + 4 × 4096 + 16 × 128 + 16 × 4096 = 100 352 bit = 3136 word × 32 bit = 12,25 KiB.

\begin{notebox}[Quyết định baseline]
\textbf{D-28 (kích thước và độ trễ).} Mỗi bank: A 256 word, B 256 word, C 1024 word, BIAS 32 word, tổng 1568 word mỗi bank, 3136 word cho cả hai bank (12,25 KiB). Mỗi sub-bank là SRAM một cổng. Độ trễ đọc đúng 1 chu kỳ; ghi có hiệu lực ở cạnh clock cuối chu kỳ lệnh; không có backpressure, \sig{o_dw_ready} = \sig{o_dr_ready} = 1.
\end{notebox}

\subsubsection{Truy cập phía tính}

\paragraph{Cổng phía tính (b = \sigc{i_comp_bank_*})}
\begin{itemize}
\item Đọc A ($i_0$, k): sub-bank r (0 ≤ r < 8) của A[b] đọc địa chỉ \{$i_0$[4:3], k[4:2]\}. Lane r = byte k[1:0] của sub-bank r; k[1:0] được chốt để chọn byte ở chu kỳ sau
\item Đọc B (k, $j_0$): sub-bank s (s = 0, 1) của B[b] đọc \{k, $j_0$[4:3]\}. Lane c = byte (c mod 4) của sub-bank c/4
\item Đọc BIAS ($j_0$): sub-bank c (0 ≤ c < 8) của BIAS[b] đọc $j_0$[4:3]; lane c = sub-bank c
\item Ghi C INT32 (i, $j_0$): lane c ghi sub-bank c của C[b] tại \{i[4:0], $j_0$[4:3]\}, dữ liệu \sigc{i_c_wr_data[32c +: 32]}, cả 4 byte, nếu \sig{mask[c]}
\item Ghi C INT8 (i, $j_0$): word w (w = 0, 1) của đoạn hàng có địa chỉ a = \{i, $j_0$/4 + w\}, nằm ở sub-bank s = \{$j_0$[4:3], w\}, địa chỉ cục bộ a » 3 = i (chiếm 32 word đầu của sub-bank). Sub-bank s nhận word w = s[0] nếu s[2:1] = $j_0$[4:3]: dữ liệu \sigc{i_c_wr_data[32w +: 32]}, byte enable e = \sig{mask[4w+e]}. Các sub-bank còn lại không được chạm
\end{itemize}

\subsubsection{Truy cập phía DMA}

\begin{spectable}{|L{0.14}|L{0.22}|L{0.3}|L{0.34}|}{Ánh xạ cổng DMA (a = địa chỉ word logic)}{tab:5-3}{\textbf{Cổng} & \textbf{Buffer} & \textbf{Sub-bank} & \textbf{Địa chỉ cục bộ}}
Ghi \sig{i_dw_*} & A, \sig{i_dw_buf} = 0 & a[5:3] (= i[2:0]) & \{a[7:6], a[2:0]\} \\\hline
Ghi \sig{i_dw_*} & B, \sig{i_dw_buf} = 1 & a[0] (= jw[0]) & a[7:1] \\\hline
Ghi \sig{i_dw_*} & BIAS, \sig{i_dw_buf} = 2 & a[2:0] & a[4:3] \\\hline
Đọc \sig{i_dr_*} & C & a[2:0] & a[9:3] \\\hline

\end{spectable}

\begin{itemize}
\item Ghi \sig{i_dw_*}: buffer \sig{i_dw_buf}, bank \sig{i_dw_bank}, cả 4 byte của \sig{i_dw_data}; \sig{i_dw_buf} = 3 là không hợp lệ (AS-01)
\item Đọc \sig{i_dr_*}: C[\sig{i_dr_bank}], sub-bank và địa chỉ cục bộ theo bảng. Số sub-bank và bank được chốt để chọn dữ liệu ra ở chu kỳ sau
\end{itemize}

\begin{notebox}[Vì sao không bao giờ xung đột]
$i_0$ và $j_0$ là bội của 8, nên 8 hàng của một đọc A rơi vào 8 sub-bank khác nhau, 8 cột của một đọc BIAS hoặc ghi C INT32 cũng vậy; hai word của một đọc B hoặc ghi C INT8 nằm liền nhau nên rơi vào 2 sub-bank khác nhau. Ở INT8, \{i, jw\} = 8i + jw và 8i chia hết cho 8, nên sub-bank chỉ phụ thuộc jw. Mô hình kiểm đủ mọi ($i_0$, k), (k, $j_0$) và mọi đoạn hàng C.
\end{notebox}

\subsubsection{Chọn bên truy cập}

\paragraph{Mux cổng của mỗi sub-bank (X, b, s)}
\begin{itemize}
\item \sig{dma_sel} = lệnh DMA (ghi \sig{i_dw_*} cho A, B, BIAS; đọc \sig{i_dr_*} cho C) có buffer X, bank b và sub-bank s
\item \sig{cmp_sel} = lệnh phía tính tương ứng (\sig{i_a_rd_en}, \sig{i_b_rd_en}, \sig{i_bias_rd_en}, \sig{i_c_wr_en} với byte enable khác 0 cho s) và \sig{i_comp_bank_X} = b
\item RAM: \sigc{en = dma_sel | cmp_sel}; địa chỉ, dữ liệu, \sig{we}, byte enable lấy từ bên được chọn. Nếu cả hai cùng chọn, DMA thắng
\item Không cần tín hiệu pha: quyền truy cập do bank và lịch của FSM bảo đảm
\end{itemize}

\begin{notebox}[Quyết định baseline]
\textbf{D-12 (trọng tài cổng SRAM).} \sig{en} của mỗi sub-bank RAM = \sig{dma_sel} | \sig{cmp_sel}; khi trùng, DMA thắng. Đồng thời assertion AS-02 (chỉ mô phỏng) khẳng định tình huống này không bao giờ xảy ra khi \sig{i_comp_active} = 1, vì FSM đã cấp bank khác nhau cho hai phía. Lý do: hành vi phải xác định được trong RTL; assertion là cơ chế phát hiện lỗi lập lịch.
\end{notebox}

% ===================== 6 =====================
\subsection{Module con}\label{sec:6}

\begin{spectable}{|L{0.26}|L{0.15}|L{0.14}|L{0.45}|}{Tín hiệu giao tiếp của \sigc{gemm_spm_ram} (tham số DEPTH, BE)}{tab:6-1}{\textbf{Tên tín hiệu} & \textbf{Độ rộng} & \textbf{Hướng} & \textbf{Mô tả}}
\sig{clk} & 1 & Input & Clock của khối.\newline Các thanh ghi cập nhật tại cạnh lên. \\\hline

\sig{i_en} & 1 & Input & Cho phép truy cập RAM tại cạnh lên.\newline Khi bằng 0, không có lệnh đọc hoặc ghi mới. \\\hline
\sig{i_we} & 1 & Input & Chọn thao tác RAM khi i\_en = 1.\newline Giá trị 1: ghi; 0: đọc đồng bộ. \\\hline

\sig{i_be} & 4 & Input & Cho phép ghi từng byte của word.\newline Khi tham số BE = 1, bit b điều khiển byte b; khi BE = 0, bỏ qua mặt nạ và ghi cả word. \\\hline

\sig{i_addr} & $\lceil\log_2 \text{DEPTH}\rceil$ & Input & Địa chỉ word trong RAM con, từ 0 đến DEPTH − 1.\newline Được lấy mẫu cùng i\_en tại cạnh lên. \\\hline

\sig{i_wdata} & 32 & Input & Word 32 bit ghi vào RAM khi i\_en và i\_we cùng bằng 1.\newline Các byte thực sự được cập nhật theo mặt nạ i\_be. \\\hline
\sig{o_rdata} & 32 & Output & Word 32 bit đọc từ RAM.\newline Hợp lệ một chu kỳ sau lệnh có i\_en = 1 và i\_we = 0. \\\hline

\end{spectable}

\begin{spectable}{|L{0.2}|L{0.8}|}{Module con theo buffer}{tab:6-2}{\textbf{Module} & \textbf{Nội dung}}
\sig{gemm_spm_a} & 16 RAM (NB × 8), mỗi RAM 32 × 32 bit; phục vụ đọc vector A và ghi word DMA; chốt k[1:0] và bank để chọn byte trả về. \\\hline

\sig{gemm_spm_b} & 4 RAM (NB × Q), mỗi RAM 128 × 32 bit; phục vụ đọc vector B và ghi word DMA; chốt bank để chọn dữ liệu trả về. \\\hline

\sig{gemm_spm_bias} & 16 RAM (NB × 8), mỗi RAM 4 × 32 bit; hiện thực bằng register file vì mỗi RAM chỉ sâu 4 word, đọc vẫn đồng bộ; đọc vector bias và ghi word DMA. \\\hline

\sig{gemm_spm_c} & 16 RAM (NB × 8), mỗi RAM 128 × 32 bit có byte enable; ghi đoạn hàng INT32/INT8 và đọc word DMA; ghép word theo $j_0$ khi ghi INT8, chốt bank và sub-bank để chọn dữ liệu đọc. \\\hline

\end{spectable}

Bốn module con theo buffer chỉ chứa RAM và thanh ghi dữ liệu nên không có \sig{rst_n}. Việc giải mã buffer (\sig{i_dw_buf}) và thanh ghi \sig{o_dr_rvalid} nằm ở \sig{gemm_local_sram}.

\begin{spectable}{|L{0.26}|L{0.15}|L{0.14}|L{0.45}|}{Tín hiệu giao tiếp của \sigc{gemm_spm_a}}{tab:6-3}{\textbf{Tên tín hiệu} & \textbf{Độ rộng} & \textbf{Hướng} & \textbf{Mô tả}}
\sig{clk} & 1 & Input & Clock của khối. \\\hline
\sig{i_dw_valid} & 1 & Input & Ghi một word DMA vào A; bằng \sig{i_dw_valid} ∧ (\sig{i_dw_buf} = 0). \\\hline
\sig{i_dw_bank} & 1 & Input & Bank đích của word ghi. \\\hline
\sig{i_dw_addr} & 8 & Input & Địa chỉ word logic \{i[4:0], kw[2:0]\}, 0 … 255. \\\hline
\sig{i_dw_data} & 32 & Input & Word ghi, cả 4 byte. \\\hline
\sig{i_comp_bank} & 1 & Input & Bank phía tính, bằng \sig{i_comp_bank_a}. \\\hline
\sig{i_rd_en} & 1 & Input & Cho phép đọc vector A, bằng \sig{i_a_rd_en}. \\\hline
\sig{i_rd_row} & 5 & Input & Hàng đầu $i_0$, bội của 8. \\\hline
\sig{i_rd_k} & 5 & Input & Chỉ số k. \\\hline
\sig{o_rd_data} & 64 & Output & Vector 8 byte, lane r tại [8r +: 8]; hợp lệ 1 chu kỳ sau \sig{i_rd_en}. \\\hline

\end{spectable}

\begin{spectable}{|L{0.26}|L{0.15}|L{0.14}|L{0.45}|}{Tín hiệu giao tiếp của \sigc{gemm_spm_b}}{tab:6-4}{\textbf{Tên tín hiệu} & \textbf{Độ rộng} & \textbf{Hướng} & \textbf{Mô tả}}
\sig{clk} & 1 & Input & Clock của khối. \\\hline
\sig{i_dw_valid} & 1 & Input & Ghi một word DMA vào B; bằng \sig{i_dw_valid} ∧ (\sig{i_dw_buf} = 1). \\\hline
\sig{i_dw_bank} & 1 & Input & Bank đích của word ghi. \\\hline
\sig{i_dw_addr} & 8 & Input & Địa chỉ word logic \{k[4:0], jw[2:0]\}, 0 … 255. \\\hline
\sig{i_dw_data} & 32 & Input & Word ghi, cả 4 byte. \\\hline
\sig{i_comp_bank} & 1 & Input & Bank phía tính, bằng \sig{i_comp_bank_b}. \\\hline
\sig{i_rd_en} & 1 & Input & Cho phép đọc vector B, bằng \sig{i_b_rd_en}. \\\hline
\sig{i_rd_k} & 5 & Input & Chỉ số k của hàng B. \\\hline
\sig{i_rd_col} & 5 & Input & Cột đầu $j_0$, bội của 8. \\\hline
\sig{o_rd_data} & 64 & Output & Vector 8 byte, lane c tại [8c +: 8]; hợp lệ 1 chu kỳ sau \sig{i_rd_en}. \\\hline

\end{spectable}

\begin{spectable}{|L{0.26}|L{0.15}|L{0.14}|L{0.45}|}{Tín hiệu giao tiếp của \sigc{gemm_spm_bias}}{tab:6-5}{\textbf{Tên tín hiệu} & \textbf{Độ rộng} & \textbf{Hướng} & \textbf{Mô tả}}
\sig{clk} & 1 & Input & Clock của khối. \\\hline
\sig{i_dw_valid} & 1 & Input & Ghi một word DMA vào BIAS; bằng \sig{i_dw_valid} ∧ (\sig{i_dw_buf} = 2). \\\hline
\sig{i_dw_bank} & 1 & Input & Bank đích của word ghi. \\\hline
\sig{i_dw_addr} & 5 & Input & Chỉ số j của bias, 0 … 31. \\\hline
\sig{i_dw_data} & 32 & Input & Word bias INT32. \\\hline
\sig{i_comp_bank} & 1 & Input & Bank phía tính, bằng \sig{i_comp_bank_bias}. \\\hline
\sig{i_rd_en} & 1 & Input & Cho phép đọc vector bias, bằng \sig{i_bias_rd_en}. \\\hline
\sig{i_rd_col} & 5 & Input & Cột đầu $j_0$, bội của 8. \\\hline
\sig{o_rd_data} & 256 & Output & 8 bias INT32, lane c tại [32c +: 32]; hợp lệ 1 chu kỳ sau \sig{i_rd_en}. \\\hline

\end{spectable}

\begin{spectable}{|L{0.26}|L{0.15}|L{0.14}|L{0.45}|}{Tín hiệu giao tiếp của \sigc{gemm_spm_c}}{tab:6-6}{\textbf{Tên tín hiệu} & \textbf{Độ rộng} & \textbf{Hướng} & \textbf{Mô tả}}
\sig{clk} & 1 & Input & Clock của khối. \\\hline
\sig{i_dr_valid} & 1 & Input & Yêu cầu đọc một word DMA từ C. \\\hline
\sig{i_dr_bank} & 1 & Input & Bank nguồn của word đọc. \\\hline
\sig{i_dr_addr} & 10 & Input & Địa chỉ word logic trong C. \\\hline
\sig{o_dr_rdata} & 32 & Output & Word đọc, hợp lệ 1 chu kỳ sau \sig{i_dr_valid}. \\\hline
\sig{i_comp_bank} & 1 & Input & Bank phía tính, bằng \sig{i_comp_bank_c}. \\\hline
\sig{i_wr_en} & 1 & Input & Ghi một đoạn hàng, bằng \sig{i_c_wr_en}. \\\hline
\sig{i_wr_fmt} & 1 & Input & 0: INT32; 1: INT8. \\\hline
\sig{i_wr_row} & 5 & Input & Chỉ số hàng i. \\\hline
\sig{i_wr_col} & 5 & Input & Cột đầu $j_0$, bội của 8. \\\hline
\sig{i_wr_data} & 256 & Input & Dữ liệu đoạn hàng, bố cục như \sig{i_c_wr_data}. \\\hline
\sig{i_wr_mask} & 8 & Input & Mặt nạ phần tử, như \sig{i_c_wr_mask}. \\\hline

\end{spectable}

\begin{Verbatim}[frame=leftline,framerule=2pt,rulecolor=\color{brand},xleftmargin=6pt,samepage=true]
// gemm_spm_ram: RAM một cổng, đọc đồng bộ (tham khảo)
logic [31:0] mem [DEPTH];
always_ff @(posedge clk)
  if (i_en) begin
    if (i_we) begin
      for (int e = 0; e < 4; e++)
        if (!BE || i_be[e]) mem[i_addr][8*e +: 8] <= i_wdata[8*e +: 8];
    end else
      o_rdata <= mem[i_addr];
  end
\end{Verbatim}

% ===================== 7 =====================
\subsection{Timing}\label{sec:7}
\begin{spectable}{|L{0.42}|L{0.29}|L{0.29}|}{Thời điểm (t = chu kỳ có lệnh)}{tab:7-1}{\textbf{Cổng} & \textbf{Lệnh} & \textbf{Kết quả}}
A, B, BIAS đọc & t & dữ liệu ở t + 1 \\\hline

C ghi (VPU), A/B/BIAS ghi (DMA) & t & có trong RAM từ t + 1 \\\hline

C đọc (DMA) & t & \sig{o_dr_rvalid} = 1, dữ liệu ở t + 1 \\\hline

\end{spectable}

\begin{figure}[H]\centering
\begin{wave}{7}{-13.9}
\sname{0}{i\_comp\_active} \hi{0}{7}{-0.2}
\sname{-1.3}{i\_comp\_bank\_a} \bus{0}{7}{-1.3}{0}
\sname{-2.6}{i\_a\_rd\_en} \hi{0}{2}{-2.8} \lo{2}{7}{-2.8}
\sname{-3.9}{i\_a\_rd\_k} \bus{0}{1}{-3.9}{k0} \bus{1}{2}{-3.9}{k1}
\sname{-5.2}{o\_a\_rd\_data} \bus{1}{2}{-5.2}{A(k0)} \bus{2}{3}{-5.2}{A(k1)}
\sname{-6.5}{i\_dw\_valid} \hi{0}{3}{-6.7} \lo{3}{7}{-6.7}
\sname{-7.8}{i\_dw\_bank} \bus{0}{3}{-7.8}{1}
\sname{-9.1}{i\_dw\_addr} \bus{0}{1}{-9.1}{a0} \bus{1}{2}{-9.1}{a1} \bus{2}{3}{-9.1}{a2}
\sname{-10.4}{i\_dr\_valid} \lo{0}{4}{-10.6} \hi{4}{5}{-10.6} \lo{5}{7}{-10.6}
\sname{-11.7}{o\_dr\_rvalid} \lo{0}{5}{-11.9} \hi{5}{6}{-11.9} \lo{6}{7}{-11.9}
\sname{-13.0}{o\_dr\_rdata} \bus{5}{6}{-13.0}{C[a]}
\end{wave}
\caption{Độ trễ đọc cố định 1 chu kỳ và đồng thời hai bank. Phía tính đọc A bank 0 ở chu kỳ 0 và 1 (\sigc{i_comp_bank_a} = 0), dữ liệu A(k0), A(k1) ở chu kỳ 1 và 2. Trong cùng các chu kỳ đó DMA ghi ba word vào A bank 1 (\sigc{i_dw_bank} = 1), không tranh sub-bank. Ở chu kỳ 4 DMA đọc C bank 1 (phía tính dùng bank C 0); \sigc{o_dr_rvalid} và dữ liệu xuất hiện đúng ở chu kỳ 5. Mọi cổng không có backpressure}\label{fig:7-1}
\end{figure}

\begin{itemize}
\item Ghi tại chu kỳ t có hiệu lực ở cạnh clock cuối chu kỳ t; đọc cùng địa chỉ ở chu kỳ t + 1 thấy dữ liệu mới
\item Đọc và ghi cùng sub-bank trong cùng một chu kỳ bị cấm (cổng RAM một cổng); thiết kế chỉ cho phép khi hai truy cập thuộc hai bank khác nhau
\end{itemize}

% ===================== 8 =====================
\subsection{Reset, hiện thực, kiểm chứng}\label{sec:8}

\paragraph{Reset}
\begin{itemize}
\item \sig{rst_n} = 0 hoặc \sig{i_soft_clr}: \sig{o_dr_rvalid} về 0. Thanh ghi chọn byte, bank, sub-bank và nội dung RAM không bị reset
\item \sig{o_dr_rvalid} là flip-flop trạng thái duy nhất của module: \sigc{always_ff @(posedge clk or negedge rst_n)}, nhánh reset và nhánh \sig{i_soft_clr} đều về 0, ngược lại nhận \sig{i_dr_valid}
\end{itemize}

\begin{notebox}[Quyết định baseline]
\textbf{D-21 (kiểu always\_ff).} Flip-flop trạng thái (\sig{o_dr_rvalid}) dùng \sigc{always_ff @(posedge clk or negedge rst_n)}; thanh ghi chỉ chứa dữ liệu (thanh ghi chọn byte, bank, sub-bank; RAM; thanh ghi đầu ra RAM) dùng \sigc{always_ff @(posedge clk)} và không có giá trị reset. \sig{soft_clr} là xóa đồng bộ, không đi vào chân reset.
\end{notebox}

\paragraph{Hiện thực}
\begin{itemize}
\item RAM viết theo mẫu \sig{gemm_spm_ram} để công cụ suy ra block RAM hoặc LUTRAM có đọc đồng bộ. Không được dùng RAM đọc bất đồng bộ cho A, B, C, BIAS
\item Không có mux tổ hợp từ \sig{*_rd_en} tới dữ liệu ra; đường ra là: thanh ghi đầu ra RAM → mux bank → mux byte/lane. Thanh ghi chọn bank, byte (k[1:0]) và sub-bank được nạp ở mỗi chu kỳ từ các cổng địa chỉ
\item Mỗi buffer trong file riêng; NB và số sub-bank dựng bằng \sig{generate}
\end{itemize}

\paragraph{Assertion bắt buộc (SVA)}
\begin{itemize}
\item \textbf{AS-01}: \sig{i_dw_valid} → \sig{i_dw_buf} ≠ 3 ∧ địa chỉ trong miền của buffer (A, B < 256; BIAS < 32)
\item \textbf{AS-02}: khi \sig{i_comp_active} = 1, không sub-bank nào có \sig{dma_sel} ∧ \sig{cmp_sel} cùng chu kỳ; ở mức bank: \sig{i_dw_valid} ∧ \sig{i_dw_buf} ∈ \{A, B\} → \sig{i_dw_bank} ≠ \sig{i_comp_bank_X} với X = \sig{i_dw_buf}; \sig{i_dw_valid} ∧ \sig{i_dw_buf} = BIAS ∧ \sig{i_bias_rd_en} → \sig{i_dw_bank} ≠ \sig{i_comp_bank_bias}; \sig{i_dr_valid} → \sig{i_dr_bank} ≠ \sig{i_comp_bank_c}. Thành phần BIAS chỉ áp khi phía tính đang đọc bias, vì khi \sig{POST_EN} = 0 Main FSM được phép nạp BIAS cho job sau vào bank trùng \sig{i_comp_bank_bias} (không phải xung đột). Ràng buộc lập lịch đầy đủ cho BIAS nằm ở AS-04 của Main FSM. Đây là assertion của quyết định D-12: RTL vẫn xác định (DMA thắng) nếu vi phạm, assertion phát hiện lỗi lập lịch
\item \textbf{AS-03}: \sig{i_a_rd_en} ∨ \sig{i_b_rd_en} ∨ \sig{i_bias_rd_en} ∨ \sig{i_c_wr_en} → \sig{i_comp_active}
\item \textbf{AS-04}: \sig{i_a_rd_row}, \sig{i_b_rd_col}, \sig{i_bias_rd_col}, \sig{i_c_wr_col} là bội của 8 khi có lệnh
\item \textbf{AS-05}: \sig{o_dw_ready} = \sig{o_dr_ready} = 1 ở mọi chu kỳ; \sig{o_dr_rvalid} bằng \sig{i_dr_valid} của chu kỳ trước
\item \textbf{AS-06}: \sig{i_c_wr_en} ∧ \sig{i_c_wr_fmt} = 1 → \sig{i_c_wr_data[255:64]} = 0
\end{itemize}

\paragraph{Kiểm tra chức năng}
\begin{itemize}
\item Ghi A, B, BIAS qua cổng DMA rồi đọc qua cổng array/VPU: mọi (i$_0$, k), (k, $j_0$), $j_0$ khớp giá trị đã ghi, đúng 1 chu kỳ sau \sig{*_rd_en}
\item Ghi C INT32 và INT8 qua cổng VPU với mask ngẫu nhiên rồi đọc qua cổng DMA: byte/word có mask = 0 giữ nguyên giá trị cũ, phần tử có mask = 1 khớp
\item Phía tính và DMA truy cập đồng thời hai bank khác nhau của cùng buffer: không sai lệch dữ liệu ở cả hai phía
\item Số sub-bank bị chạm trong mỗi truy cập song song bằng đúng bảng ở mục điều kiện phục vụ đồng thời, không sub-bank nào bị chạm hai lần
\end{itemize}

\Needspace{8\baselineskip}
\begin{notebox}[Kiểm chứng bằng mô hình]
(Mô hình nằm ở \sig{gemm-accel/model/}, đã chạy lại với baseline ngày 12/11/2026; kết quả đạt.) \sig{gemm_sram_model.py}: ánh xạ sub-bank và địa chỉ cục bộ là song ánh với P = 8; 240 cấu hình ngẫu nhiên (M, N, K, bank, NB = 2) nạp A, B, BIAS qua cổng DMA rồi đọc qua cổng array/VPU, ghi C INT32 và INT8 qua cổng VPU rồi đọc qua cổng DMA: mọi giá trị khớp và không truy cập song song nào chạm hai lần một sub-bank.
\end{notebox}

\end{document}
